\thispagestyle{fancy}
\fancyhead[R]{A.U 2025-2026}
\fancyhead[L]{Conclusion générale et perspectives}

\chapter*{Conclusion générale et perspectives}
\addcontentsline{toc}{chapter}{Conclusion générale et perspectives}

Le travail présenté dans ce rapport s'inscrit dans le cadre du projet de fin d'études
réalisé au sein de la société \textbf{CLEDISS}, éditeur tunisien de solutions logicielles.
Il répond à une problématique concrète et quotidienne de l'équipe technique :
l'absence d'un outil unifié permettant de superviser l'état des serveurs hébergeant les
applications de production et d'y intervenir rapidement en cas d'incident.

Avant ce projet, la détection d'une anomalie reposait essentiellement sur le signalement
des utilisateurs finaux, et chaque intervention imposait une connexion manuelle au
serveur concerné. Cette approche réactive engendrait des délais de résolution importants
et ne laissait aucune trace exploitable des actions réalisées.

\section*{Bilan du travail réalisé}
\addcontentsline{toc}{section}{Bilan du travail réalisé}

La solution développée est une \textbf{plateforme de supervision et d'administration
d'infrastructure} constituée de quatre composants complémentaires : un agent de collecte
Python déployé sur chaque serveur supervisé, une interface d'administration web
développée en React.js, une application mobile React Native, et un backend Node.js
exposant une API REST commune aux deux clients.

Conformément à la méthodologie Scrum retenue, le développement s'est déroulé en un sprint
préparatoire suivi de quatre sprints de réalisation :

\begin{itemize}[label=\ding{43}]
  \item le \textbf{Sprint 0} a permis de cadrer le projet : étude de l'existant, analyse
        des besoins, choix de l'architecture N-Tiers et conception des modèles de
        données ;
  \item le \textbf{Sprint 1} a établi le socle de la plateforme : authentification
        sécurisée par jeton JWT, agent de collecte des métriques système, enregistrement
        automatique des serveurs et tableau de bord temps réel alimenté par WebSocket ;
  \item le \textbf{Sprint 2} a transformé l'outil de visualisation en outil
        d'administration : alertes automatiques par seuils avec notification par courrier
        électronique, détection dynamique des services, actions à distance sécurisées via
        SSH, journal d'audit, sauvegardes automatisées et application mobile ;
  \item le \textbf{Sprint 3} a industrialisé la chaîne de livraison : conteneurisation
        des trois composants avec Docker, \textit{pipeline} d'intégration continue en
        sept étapes sous Jenkins, analyse de la qualité du code avec SonarQube et analyse
        des vulnérabilités des images avec Trivy ;
  \item le \textbf{Sprint 4} a automatisé le déploiement : orchestration sur un cluster
        Kubernetes k3s, déploiement continu selon l'approche GitOps avec Argo CD,
        externalisation des secrets dans HashiCorp Vault, cloisonnement des droits par
        RBAC et supervision de l'infrastructure avec Prometheus et Grafana.
\end{itemize}

À l'issue de ces cinq sprints, l'ensemble des objectifs fixés au démarrage du projet a
été atteint. La plateforme est déployée en production sur le serveur VPS de l'entreprise
et supervise effectivement plusieurs serveurs, tandis que la chaîne DevSecOps assure de
bout en bout la construction, la validation et le déploiement de chaque nouvelle version
sans intervention manuelle.

\section*{Apports du projet}
\addcontentsline{toc}{section}{Apports du projet}

\subsection*{Pour l'entreprise}
La plateforme apporte à CLEDISS un gain opérationnel mesurable. La détection des
anomalies est devenue \textbf{proactive} : l'équipe technique est alertée dès le
dépassement d'un seuil, avant que le service ne se dégrade pour les utilisateurs. Les
interventions courantes — redémarrage d'un service, consultation de l'état d'un serveur —
s'effectuent désormais depuis une interface unique, en quelques secondes, et depuis un
téléphone si nécessaire. Enfin, chaque action est tracée dans un journal d'audit, ce qui
répond aux exigences de traçabilité auxquelles est soumise l'entreprise.

\subsection*{Sur le plan personnel}
Ce projet a été l'occasion de mettre en pratique, sur un cas réel et complet, les
compétences acquises durant le cursus d'ingénieur. Il m'a permis de maîtriser
l'intégralité d'une chaîne applicative moderne, du développement à l'exploitation :
conception d'une architecture N-Tiers, développement d'une API REST sécurisée,
réalisation d'interfaces web et mobiles, et surtout appropriation concrète de la culture
\textbf{DevSecOps}, où la sécurité et la qualité sont intégrées au processus de
construction plutôt que contrôlées après coup.

L'immersion en entreprise a également apporté une dimension humaine essentielle :
échanges réguliers avec l'encadrant technique, présentation des livrables à la fin de
chaque sprint, prise en compte des retours et arbitrages entre l'idéal technique et les
contraintes de délai.

\section*{Difficultés rencontrées}
\addcontentsline{toc}{section}{Difficultés rencontrées}

Plusieurs difficultés techniques ont jalonné le projet et ont donné lieu à des solutions
qui ont notablement enrichi la conception :

\begin{itemize}[label=\ding{43}]
  \item \textbf{L'hétérogénéité du parc supervisé.} Imposer une liste figée de services
        aurait restreint la plateforme à un ensemble prédéfini. La solution retenue —
        une détection dynamique des unités actives par l'agent — s'est révélée à la fois
        plus souple et plus sûre, la liste détectée servant également de liste blanche
        avant toute exécution de commande.
  \item \textbf{La sécurisation des actions à distance.} Exécuter une commande
        privilégiée sur un serveur depuis une interface web constitue une opération
        sensible. Un dispositif à plusieurs niveaux a été mis en place : double
        vérification d'autorisation, validation du nom de service par expression
        régulière, journalisation systématique et notification par courrier électronique.
  \item \textbf{La contrainte de ressources du serveur.} Le VPS de l'entreprise ne
        pouvait accueillir une distribution Kubernetes complète. Le choix de k3s a permis
        de conserver l'intégralité de l'interface de programmation de Kubernetes tout en
        respectant cette contrainte matérielle.
  \item \textbf{La gestion des secrets.} Le versionnement des manifestes de déploiement
        exposait initialement les identifiants applicatifs dans le dépôt Git. Leur
        externalisation dans HashiCorp Vault a corrigé cette faiblesse et constitue l'un
        des apports majeurs du dernier sprint.
\end{itemize}

\section*{Perspectives d'évolution}
\addcontentsline{toc}{section}{Perspectives d'évolution}

La plateforme est fonctionnelle et déployée, mais plusieurs axes d'amélioration ont été
identifiés et pourraient faire l'objet de développements ultérieurs :

\begin{itemize}[label=\ding{220}]
  \item \textbf{Notifications multicanal.} Compléter les notifications par courrier
        électronique par des canaux à latence réduite tels que Slack, Microsoft Teams ou
        les notifications \textit{push} de l'application mobile.
  \item \textbf{Seuils dynamiques et détection d'anomalies.} Remplacer les seuils fixes
        par des seuils calculés à partir du comportement historique de chaque serveur,
        afin de détecter les dérives progressives que des seuils absolus ne révèlent pas.
  \item \textbf{Authentification par clé SSH.} Substituer aux mots de passe SSH une
        authentification par clés, gérées dans Vault, ce qui renforcerait encore la
        sécurité des actions à distance.
  \item \textbf{Gestion multi-utilisateurs.} Introduire une gestion fine des rôles au
        niveau applicatif afin de distinguer les profils consultation, exploitation et
        administration, et de restreindre les actions sensibles à certains opérateurs.
  \item \textbf{Haute disponibilité.} Étendre le cluster à plusieurs n\oe uds et
        augmenter le nombre de réplicas des pods applicatifs afin de supprimer le point
        de défaillance unique que constitue aujourd'hui le serveur unique.
  \item \textbf{Supervision applicative approfondie.} Étendre la collecte au-delà des
        métriques système en y intégrant les journaux applicatifs et les temps de réponse
        des services métiers, pour une vision complète de la santé des applications.
\end{itemize}

En définitive, ce projet de fin d'études a permis de livrer à CLEDISS une solution
opérationnelle répondant à un besoin réel, tout en constituant une expérience formatrice
particulièrement complète, à la croisée du développement logiciel, de l'administration
système et de l'ingénierie DevSecOps.
